\documentclass{article}
\usepackage{amsmath}
\begin{document}
\section{Figures}
This document has a figure, which is a float:
TeX moves it to the top of a page or to a page of its own,
so the source lines of the figure end up far from the surrounding text.

\begin{figure}[t]
  \centering
  \rule{0.6\textwidth}{4cm}
  \caption{A black rectangle stands in for a picture.
    The caption spans two source lines.}
  \label{fig:rect}
\end{figure}

Figure~\ref{fig:rect} shows a rectangle.
The text after the figure environment continues on the page where the float was defined,
which may be before or after the float itself.

\begin{figure}[p]
  \centering
  \fbox{\parbox{0.5\textwidth}{A framed paragraph inside a float that goes on a page of floats.
  It has two lines in the source.}}
  \caption{A float page.}
\end{figure}

\section{After the Figures}
Some mathematics in display:
\[
  e^{i\pi} + 1 = 0
\]
and a multi-line display:
\begin{gather}
  a = b + c \\
  d = e + f
\end{gather}
A minipage pair side by side:

\noindent
\begin{minipage}{0.45\textwidth}
Left minipage text, which wraps inside its own narrow box.
\end{minipage}\hfill
\begin{minipage}{0.45\textwidth}
Right minipage text, which also wraps inside a narrow box.
\end{minipage}

\vspace{1em}
A tabular with a rule:
\begin{center}
\begin{tabular}{|l|l|}
  \hline
  key & value \\
  \hline
  one & 1 \\
  two & 2 \\
  \hline
\end{tabular}
\end{center}
The end of the figure document.
\end{document}
